% =============================================================
%  Conclusion générale et perspectives
% =============================================================
\chapter*{Conclusion générale et perspectives}
\addcontentsline{toc}{chapter}{Conclusion générale et perspectives}
\markboth{Conclusion générale et perspectives}
         {Conclusion générale et perspectives}

Ce projet de fin d'études, réalisé au sein de \textbf{Sagemcom Ezzahra},
a abouti à la conception, au développement et au déploiement d'une
solution complète de tests Telnet automatisés pour les équipements réseau
embarqués. On est parti d'un processus manuel lent et fragmenté pour
livrer une plateforme web fullstack, intégrée dans une infrastructure
DevOps fonctionnelle et déployée automatiquement.

Sur le \textbf{plan applicatif}, \textbf{Telnet Test Manager} change
concrètement la façon dont les firmwares sont validés. L'application
est structurée selon les principes de la \textit{Clean Architecture}
en couches découplées (Domain, Application, Infrastructure, Interfaces),
ce qui facilite la maintenance et les évolutions futures. Le moteur
Telnet, basé sur les \textbf{Worker Threads} de Node.js, prend en
charge quatre modes~: commande unitaire, séquence pas-à-pas, supervision
continue et tests multi-slots en parallèle. Les résultats sont diffusés
en temps réel via \textbf{WebSocket} et l'interface \textbf{React~18/TypeScript}
est disponible en trois langues. Le contrôle d'accès \gls{rbac} à trois
rôles --- Administrateur, Ingénieur et Technicien --- gère les permissions
de façon claire et sécurisée.

Sur le \textbf{plan DevOps}, l'application est conteneurisée avec
\textbf{Docker} (build multi-stage qui ramène l'image frontend de
$\sim$900~Mo à $\sim$25~Mo en production) et orchestrée par un chart
\textbf{Helm} v0.3.0 sur \textbf{Kubernetes}. L'autoscaling horizontal
a été validé en conditions réelles~: lors des tests de charge, le CPU
a atteint 267\,\% du seuil, déclenchant automatiquement la montée de
1 à 5 répliques. Le \textbf{pipeline GitHub Actions} en trois jobs
gère l'intégration continue, les audits de sécurité et le déploiement
en 2~minutes~33~secondes via un runner auto-hébergé. \textbf{Prometheus}
et \textbf{Grafana} supervisent les métriques en temps réel et envoient
des alertes e-mail automatiques sur dépassement du seuil CPU.

Sur le \textbf{plan des tests}, les résultats sont clairs~: 24~tests
unitaires Jest passés à 100\,\%, 14~scénarios fonctionnels couvrant
tous les besoins identifiés, sans aucun échec sur l'ensemble de la
batterie de tests.

Pour l'équipe d'ingénierie des tests de \textbf{Sagemcom Ezzahra},
les bénéfices sont directs. Remplacer les terminaux manuels par une
interface web centralisée réduit le temps des campagnes de validation
et supprime les erreurs de saisie. Le journal d'audit et les rapports
automatiques permettent de suivre chaque campagne précisément.
L'architecture scalable et le déploiement automatisé donnent à l'équipe
une base solide pour faire évoluer la plateforme.

\section*{Perspectives d'évolution}

Les fondations établies par ce projet ouvrent plusieurs axes
d'amélioration pour les versions futures~:

\begin{itemize}
  \item \textbf{Support du protocole SSH~:} étendre le moteur de
        connexion pour couvrir les équipements modernes refusant
        les connexions Telnet non chiffrées~;

  \item \textbf{Analyse intelligente des résultats~:} appliquer
        des algorithmes de détection d'anomalies sur l'historique
        des tests pour anticiper les défaillances matérielles avant
        qu'elles n'affectent la production~;

  \item \textbf{Tests end-to-end automatisés~:} intégrer Cypress
        ou Playwright dans le pipeline CI/CD pour valider
        automatiquement les scénarios utilisateur complets~;

  \item \textbf{Déploiement cloud hybride~:} migrer vers AWS EKS
        ou GCP GKE pour étendre la portée de l'infrastructure à
        plusieurs sites géographiques de Sagemcom~;

  \item \textbf{Planification des campagnes~:} implémenter un
        module de scheduling pour lancer des campagnes de tests
        automatiques sur des plages horaires programmées.
\end{itemize}

\vspace{0.5cm}
\noindent Ce projet montre qu'en combinant une architecture logicielle
solide (\textit{Clean Architecture}, WebSocket, \gls{rbac}) avec des
pratiques \gls{devops} concrètes (Docker, Kubernetes, CI/CD,
Prometheus/Grafana), il est possible de remplacer un processus manuel
fragmenté par une plateforme automatisée, traçable et scalable. Le
résultat est une infrastructure de tests qui s'adapte à la charge, se
déploie seule et laisse une trace de chaque action --- exactement ce
dont l'équipe de Sagemcom Ezzahra avait besoin.
